% ============================================================================
% RESUMEN + PALABRAS CLAVES - Según Anexo 7 de la Pauta UTEM
% ============================================================================
% El resumen debe dar cuenta en forma clara y simple del contenido de la obra.
% Extensión máxima: 1 página.
% Debe incluir:
%   a) El objetivo del trabajo
%   b) Método o procedimiento utilizado
%   c) Conclusiones o resultados obtenidos
% ============================================================================

\chapter*{RESUMEN}
\thispagestyle{fancy}

El presente trabajo aborda la ineficiencia logística y operativa en el control de acceso vehicular del Campus Ñuñoa de la Universidad Tecnológica Metropolitana (UTEM). Actualmente, el ingreso de funcionarios depende de la intervención manual del personal de conserjería, lo que genera tiempos de espera prolongados y distrae recursos humanos de labores críticas de seguridad institucional. Para resolver esta problemática, el objetivo general del proyecto es validar la viabilidad técnica y operativa de un sistema de control de acceso automatizado mediante el desarrollo de un prototipo funcional basado en Inteligencia Artificial (IA) para el reconocimiento de placas patentes chilenas.

La metodología de desarrollo se fundamenta en un enfoque híbrido que combina el marco de trabajo ágil Scrum para la gestión de las fases del proyecto, con el estándar CRISP-DM para gobernar el ciclo de vida del modelo de aprendizaje profundo. A nivel arquitectónico, se diseñó una solución bajo el paradigma de \textit{Edge Computing} (\textit{Retrofit} tecnológico), trasladando la inferencia de la red neuronal YOLOv12s directamente al nodo de borde (portería) para garantizar tiempos de respuesta inferiores a 100 ms y asegurar el funcionamiento del sistema ante caídas de red (\textit{Offline-First}).

En esta primera etapa de investigación y diseño, los resultados demuestran la viabilidad estructural de la propuesta técnica y su alto impacto institucional. El estudio de factibilidad económica confirma que la modernización del acceso, mediante la instalación de infraestructura de borde, es una inversión altamente rentable para la universidad, logrando un Valor Actual Neto (VAN) de \$5.898.595 CLP, una Tasa Interna de Retorno (TIR) del 59,47\% y un rápido Periodo de Recuperación de Capital (PRC) de 1,55 años. Este diseño teórico-práctico sienta las bases robustas para la etapa de desarrollo de software, entrenamiento del modelo empírico y pruebas de integración tecnológica correspondientes a la siguiente fase del proyecto.

\noindent\textbf{PALABRAS CLAVES:}\\
Visión Computacional, Inteligencia Artificial, Edge Computing, CRISP-DM, Evaluación de Proyectos.

\newpage
